\chapter{Codebase Progress}
\label{ch:progress}

% Boundary rule: hand-written narrative here stays stable; all
% fast-moving facts come from \texttt{generated/status.tex}, which is
% regenerated by \texttt{tools/report-status.sh} before each build and
% is gitignored. Do not hardcode module counts or SHAs in this chapter.

This chapter is a snapshot, not a history. The project's own documents
divide the timeline three ways: \texttt{docs/OVERVIEW.md} states what
works today, \texttt{docs/DELIVERED.md} records how each milestone was
earned, and \texttt{docs/ROADMAP.md} says what comes next. What follows
is the machine-counted complement to those narratives: which libraries
exist, how large each area is, and exactly which commit and toolchain
this copy of the report describes. It is regenerated on every build, so
it can never go stale the way hand-maintained counts do.

\section{Snapshot}

Build SHA: \circegitsha. Lean modules counted: \circemodulecount.

\InputIfFileExists{generated/status.tex}{%
  % generated/status.tex provides the module-area table below.
}{%
  \begin{quote}
  No generated snapshot found. Run \texttt{tools/report-status.sh} from the
  repository root (or \texttt{make pdf}, which runs it automatically).
  \end{quote}
}
